\chapter{Optimisasi Query Analitik}
\label{modul:optimisasi-query-analitik}

\section{Identitas Modul}

\begin{identitasmodul}
\textbf{Sumber materi}    & Pertemuan 7 \\
\textbf{CPMK}             & CPMK-092 \\
\textbf{Minggu pelaksanaan} & Minggu ke-9, setelah UTS \\
\textbf{Durasi tatap muka} & 120 menit \\
\textbf{Estimasi tugas}   & 4--5 jam di luar sesi \\
\textbf{Moda kerja}       & Individual \\
\textbf{Perangkat}        & PostgreSQL 17 dan DBeaver \\
\textbf{Dataset}          & NusaMart ukuran \textbf{sedang} dan Northwind sebagai pembanding \\
\textbf{Skrip pendukung}  & \berkas{sql/modul-06/} beserta \berkas{check.sql} \\
\textbf{Luaran}           & \berkas{catatan-optimisasi.md} \\
\end{identitasmodul}

\slideref{7}

\section{Capaian dan Indikator Keberhasilan}

Mahasiswa mampu membuktikan dampak partitioning, aggregate table, dan
materialized view menggunakan execution plan serta statistik buffer.
Keberhasilan bukan sekadar query lebih cepat, tetapi adanya bukti mekanisme
yang membuat perubahan tersebut bekerja.

\begin{capaian}
Pada akhir modul ini mahasiswa mampu:
\begin{enumerate}[leftmargin=1.4em]
  \item menetapkan \istilah{baseline} pengukuran yang adil sebelum melakukan
        optimisasi apa pun;
  \item memartisi tabel fakta menurut rentang bulan dan \textbf{membuktikan}
        terjadinya \istilah{partition pruning} dari keluaran
        \perintah{EXPLAIN};
  \item menjelaskan mengapa pruning gagal terjadi ketika penyaringan hanya
        dilakukan lewat tabel dimensi;
  \item membandingkan \istilah{aggregate table} dan \istilah{materialized view}
        untuk pertanyaan yang sama, dalam satuan ruang, waktu penyegaran, dan
        halaman yang dibaca;
  \item menyatakan hasil optimisasi dalam satuan halaman, bukan hanya detik,
        serta menjelaskan pengaruh cache pada pengukuran.
\end{enumerate}
\end{capaian}

Kaidah modul ini satu kalimat: \textbf{setiap tindakan optimisasi membeli
sesuatu, dan setiap pembelian ada harganya}. Jawaban ``lebih cepat'' tanpa
menyebut apa yang dibayar --- ruang, waktu pemuatan, atau kesegaran data ---
tidak memperoleh nilai penuh.

\section{Prasyarat dan Persiapan}

\subsection{Prasyarat}

\begin{itemize}
  \item Kemampuan membaca \perintah{EXPLAIN (ANALYZE, BUFFERS)} dari Modul 5,
        termasuk \perintah{shared hit} dan \perintah{shared read};
  \item pengertian page 8 KB, lebar baris, dan hubungan keduanya dengan jumlah
        halaman yang harus dibaca;
  \item pengertian indeks B-tree beserta pengaruh urutan kolomnya;
  \item skema \perintah{gudang} hasil Modul 5, lengkap dengan baris unknown dan
        constraint.
\end{itemize}

\subsection{Pre-lab}

\begin{enumerate}[leftmargin=1.4em]
  \item \textbf{Beralih ke dataset ukuran sedang.} Ini wajib, bukan pilihan.
        Pada 100 ribu baris, sequential scan memang menang dan seluruh
        pengukuran modul ini tidak akan memperlihatkan apa pun.

        Bangkitkan ulang sumber pada ukuran sedang, lalu jalankan kembali
        skrip pemuatan Anda sendiri dari Modul 2 sampai 5.
\begin{lstlisting}[style=shellgaya]
python seed/generate_nusamart.py --ukuran sedang --seed 42

# lalu jalankan berurutan skrip pemuatan Anda sendiri:
#   modul-02-star-schema/ddl.sql, load.sql
#   modul-03-scd2/ddl-scd2.sql, load-scd2.sql
#   modul-05-fisikal/ddl-fisikal.sql
\end{lstlisting}

        Inilah ujian pertama atas kaidah reproduksibilitas Modul 5: skrip yang
        benar-benar dapat dijalankan ulang dari awal akan menyelesaikan langkah
        ini tanpa penyuntingan manual. Bila skrip Anda menuntut perbaikan di
        sana-sini, catat apa saja yang perlu diubah --- itu temuan yang layak
        dilaporkan.

        Bila waktu tidak mencukupi, pulihkan
        \berkas{snapshot-modul-05-sedang.sql} yang disediakan asisten.
  \item Pulihkan basis data Northwind sebagai pembanding pada Bagian D:
\begin{lstlisting}[style=shellgaya]
./seed/unduh-northwind.sh
\end{lstlisting}
  \item Jalankan \perintah{ANALYZE} pada seluruh schema \perintah{gudang}.
        Statistik dari dataset kecil tidak berlaku untuk dataset sedang, dan
        perencana yang memakai statistik usang akan memilih rencana yang keliru.
  \item Catat ukuran awal pada \berkas{modul-06-optimisasi/baseline.md}:
\begin{lstlisting}[style=sqlgaya]
SELECT COUNT(*) AS jumlah_baris,
       pg_size_pretty(pg_relation_size('gudang.fakta_penjualan')) AS heap,
       pg_size_pretty(pg_indexes_size('gudang.fakta_penjualan'))  AS indeks
FROM   gudang.fakta_penjualan;
\end{lstlisting}
  \item Siapkan jawaban berikut pada \berkas{pre-lab.md}:
  \begin{enumerate}[label=\alph*.,leftmargin=1.4em]
    \item Tabel fakta Anda kini berisi sekitar lima juta baris. Dengan lebar
          baris hasil hitungan Modul 5, berapa halaman yang harus dibaca oleh
          satu sequential scan penuh?
    \item Bila sebuah query hanya memerlukan data satu bulan, berapa bagian
          dari halaman itu yang sebenarnya sia-sia dibaca?
    \item Menurut Anda, apa yang dibeli oleh partisi, dan apa yang dibayar?
  \end{enumerate}
\end{enumerate}

\section{Konsep Dasar}

\subsection{Partitioning dan Partition Pruning}

Partisi memecah satu tabel besar menjadi beberapa tabel fisik yang lebih kecil,
sementara bagi query ia tetap tampak sebagai satu tabel. Yang dibeli bukan
kecepatan membaca per halaman, melainkan \emph{kesempatan untuk tidak membaca
sama sekali}.

\begin{konsep}{Partition pruning}
Ketika perencana dapat membuktikan bahwa sebuah partisi mustahil memuat baris
yang dicari, partisi itu tidak dibaca sama sekali. Inilah \istilah{pruning}.

Query yang meminta data Januari pada tabel berpartisi bulanan selama tiga
tahun hanya menyentuh 1 dari 36 partisi --- membaca sekitar tiga persen
halaman dibandingkan sequential scan penuh.
\end{konsep}

Syaratnya satu, dan justru di sinilah kebanyakan kegagalan terjadi:
\textbf{penyaringan harus menyentuh kolom kunci partisi}. Perencana melakukan
pruning dengan membandingkan predikat query terhadap batas setiap partisi. Bila
predikatnya berada pada tabel lain --- misalnya
\perintah{dim\_tanggal.tanggal} --- perencana tidak memiliki dasar untuk
membuang partisi mana pun pada saat menyusun rencana.

\begin{peringatan}
Inilah kesalahan paling lazim pada partisi gudang data. Query yang menyaring
lewat \perintah{JOIN} ke dimensi tanggal tampak benar dan menghasilkan angka
yang benar, tetapi tetap membaca seluruh partisi. Penyelesaiannya bukan
membatalkan partisi, melainkan \emph{menambahkan} predikat atas kolom kunci
partisi, sehingga perencana memperoleh dasar untuk memangkas.
\end{peringatan}

Kunci partisi yang dipakai modul ini adalah \perintah{tanggal\_key}. Bentuk
\perintah{YYYYMMDD} yang dipilih pada Modul 2 kini terbayar: ia bilangan bulat,
sehingga batas partisi dapat dinyatakan sebagai rentang bilangan yang terbaca
manusia.

\subsection{Aggregate Table}

\istilah{Aggregate table} adalah tabel biasa berisi hasil agregasi yang sudah
dihitung sebelumnya. Query yang semula membaca lima juta baris fakta cukup
membaca beberapa ribu baris agregat.

Yang dibeli jelas. Yang dibayar ada tiga: ruang tambahan, waktu pemuatan
tambahan, dan --- yang paling sering dilupakan --- \textbf{kewajiban menjaga
agar isinya tetap sesuai dengan fakta}. Agregat yang tidak disegarkan setelah
fakta berubah adalah sumber angka salah yang paling sulit dilacak.

Keunggulan utamanya: penyegaran dapat dilakukan \emph{sebagian}. Bila hanya
data kemarin yang berubah, cukup hapus dan hitung ulang baris agregat untuk
kemarin.

\subsection{Materialized View}

\istilah{Materialized view} menyimpan hasil sebuah query sebagai data nyata di
disk, tetapi definisinya dikelola PostgreSQL. Kodenya jauh lebih ringkas
daripada aggregate table --- satu pernyataan \perintah{CREATE}, satu perintah
\perintah{REFRESH}.

Harganya terletak pada cara penyegarannya. \perintah{REFRESH MATERIALIZED VIEW}
menghitung ulang \emph{seluruh} isinya, betapa pun kecil bagian data yang
berubah. Pada tabel yang terus tumbuh, biaya penyegaran tumbuh bersamanya.

\begin{center}
\small
\begin{tabular}{@{}p{0.24\textwidth}>{\raggedright\arraybackslash}p{0.32\textwidth}>{\raggedright\arraybackslash}p{0.36\textwidth}@{}}
\toprule
\textbf{Aspek} & \textbf{Aggregate table} & \textbf{Materialized view} \\
\midrule
Banyaknya kode & Lebih banyak, ditulis sendiri & Sedikit, dikelola PostgreSQL \\
Penyegaran & Dapat sebagian & Selalu seluruhnya \\
Biaya penyegaran & Sebanding data yang berubah & Sebanding seluruh data \\
Penguncian saat segar & Dapat diatur sendiri & Mengunci, kecuali
  \perintah{CONCURRENTLY} \\
Risiko & Lupa menyegarkan & Penyegaran menjadi mahal seiring waktu \\
\bottomrule
\end{tabular}
\end{center}

\begin{catatan}
\perintah{REFRESH MATERIALIZED VIEW CONCURRENTLY} tidak mengunci pembaca,
tetapi memerlukan satu indeks \perintah{UNIQUE} pada view dan justru lebih
lambat daripada penyegaran biasa. Ia menukar waktu dengan ketersediaan ---
pertukaran yang masuk akal pada sistem yang dipakai sepanjang hari, dan tidak
masuk akal pada gudang data yang disegarkan tengah malam.
\end{catatan}

\subsection{Buffer, Cache, dan Pengukuran yang Adil}

Pengukuran yang tidak adil menghasilkan kesimpulan yang salah. Tiga sumber
ketidakadilan yang paling sering muncul:

\begin{enumerate}[leftmargin=1.4em]
  \item \textbf{Cache.} Eksekusi kedua atas query yang sama membaca dari
        memori. Waktunya bisa sepuluh kali lebih cepat tanpa satu pun
        perubahan pada rancangan.
  \item \textbf{Statistik usang.} Perencana yang tidak tahu tabelnya sudah
        membesar akan memilih rencana untuk tabel kecil.
  \item \textbf{Beban lain.} Laptop yang sedang menjalankan hal lain
        menghasilkan waktu yang berubah-ubah.
\end{enumerate}

\begin{konsep}{Ukuran yang tahan terhadap ketiganya}
Jumlah halaman yang disentuh --- \perintah{shared hit} ditambah
\perintah{shared read} --- tidak berubah oleh keadaan cache maupun beban
laptop. Angka inilah yang dipakai untuk membandingkan dua rancangan.

\perintah{shared read} sendiri menyatakan berapa halaman yang benar-benar
diambil dari disk, dan angka itu memang berubah menurut cache. Catat keduanya,
tetapi \emph{bandingkan} jumlahnya.
\end{konsep}

Protokol pengukuran yang dipakai modul ini: jalankan setiap query
\textbf{tiga kali}, catat eksekusi pertama dan ketiga, lalu bandingkan
berdasarkan total halaman.

\section{Studi Kasus dan Protokol Praktikum}

Empat tindakan optimisasi diuji terhadap pertanyaan analitik yang tetap.
Dataset sedang wajib digunakan agar biaya scan dan pruning terlihat.

Satu pertanyaan dipakai sepanjang Bagian A sampai C, sehingga keempat tindakan
dapat dibandingkan secara setara:

\begin{quote}
\emph{Berapa nilai penjualan per kategori produk per bulan di Provinsi Lampung
sepanjang paruh pertama 2024?}
\end{quote}

\begin{center}
\small
\begin{tabular}{@{}p{0.06\textwidth}p{0.32\textwidth}>{\raggedright\arraybackslash}p{0.54\textwidth}@{}}
\toprule
\textbf{No} & \textbf{Tindakan} & \textbf{Yang hendak dibuktikan} \\
\midrule
1 & Partisi bulanan & Pruning membuang partisi di luar rentang \\
2 & Indeks pada tabel terpartisi & Penyempitan lebih lanjut di dalam partisi \\
3 & Aggregate table & Membaca ribuan baris, bukan jutaan \\
4 & Materialized view & Hasil setara, dengan biaya penyegaran berbeda \\
\bottomrule
\end{tabular}
\end{center}

\begin{catatanasisten}
Pemuatan dataset sedang memakan waktu beberapa menit dan harus sudah selesai
sebelum sesi. Pastikan mahasiswa menjalankannya sebagai pre-lab, dan sediakan
\berkas{snapshot-modul-05-sedang.sql} bagi yang gagal. Sesi 120 menit tidak
cukup bila pemuatan lima juta baris dilakukan di kelas.
\end{catatanasisten}

\section{Alur Praktikum 120 Menit}

\begin{alurwaktu}
0--15 & Demonstrasi baseline dan cara membaca buffer. \\
15--95 & Kerja terbimbing: partisi, agregat, materialized view, dan Northwind. \\
95--110 & Checkpoint bukti pruning serta tabel sebelum--sesudah. \\
110--120 & Penjelasan analisis biaya penyegaran dan ruang. \\
\end{alurwaktu}

\section{Prosedur Praktikum Terbimbing}

\subsection{Bagian A: Menetapkan Baseline}

\langkah{A.1 Mengukur keadaan awal}{10 menit}

\begin{lstlisting}[style=sqlgaya]
EXPLAIN (ANALYZE, BUFFERS)
SELECT dp.kategori, dt.tahun_bulan, SUM(f.subtotal) AS nilai_penjualan
FROM   gudang.fakta_penjualan f
JOIN   gudang.dim_tanggal dt  ON dt.tanggal_key = f.tanggal_key
JOIN   gudang.dim_produk  dp  ON dp.produk_key  = f.produk_key
JOIN   gudang.dim_toko    dtk ON dtk.toko_key   = f.toko_key
WHERE  dtk.provinsi = 'Lampung'
  AND  dt.tanggal BETWEEN DATE '2024-01-01' AND DATE '2024-06-30'
GROUP  BY 1, 2;
\end{lstlisting}

Jalankan tiga kali. Catat pada \berkas{catatan-optimisasi.md}: jenis scan,
\perintah{shared hit}, \perintah{shared read}, jumlahnya, serta waktu eksekusi
pertama dan ketiga.

\begin{catatan}
Perhatikan selisih waktu antara eksekusi pertama dan ketiga, lalu perhatikan
bahwa jumlah halaman hampir tidak berubah. Inilah bukti langsung mengapa waktu
bukan ukuran yang dapat dipercaya untuk membandingkan rancangan --- dan
mengapa seluruh perbandingan pada modul ini memakai halaman.
\end{catatan}

\subsection{Bagian B: Menerapkan Partisi Bulanan}

\langkah{B.1 Membangun tabel berpartisi}{15 menit}

Tabel yang sudah ada tidak dapat diubah menjadi tabel berpartisi. Yang
dilakukan adalah membuat tabel berpartisi baru, memindahkan datanya, lalu
menukar namanya. Skrip \berkas{sql/modul-06/\allowbreak 01-partisi.sql}
mengerjakan langkah ini beserta pembangkitan partisi bulanan.

\begin{lstlisting}[style=sqlgaya]
CREATE TABLE gudang.fakta_penjualan_p (
  LIKE gudang.fakta_penjualan INCLUDING DEFAULTS
) PARTITION BY RANGE (tanggal_key);

-- Satu partisi per bulan, batas dinyatakan sebagai rentang YYYYMMDD.
CREATE TABLE gudang.fakta_penjualan_2024_01
  PARTITION OF gudang.fakta_penjualan_p
  FOR VALUES FROM (20240101) TO (20240201);
-- ... dan seterusnya

-- Partisi penampung untuk nilai di luar rentang mana pun.
CREATE TABLE gudang.fakta_penjualan_lain
  PARTITION OF gudang.fakta_penjualan_p DEFAULT;
\end{lstlisting}

\begin{peringatan}
Partisi \perintah{DEFAULT} wajib ada, dan alasannya berasal dari Modul 5:
baris fakta yang tanggalnya tidak diketahui menunjuk
\perintah{tanggal\_key = -1}, dan nilai itu tidak masuk rentang bulan mana
pun. Tanpa partisi penampung, pemindahan data akan gagal dengan pesan
\emph{no partition of relation found for row}.

Periksa isinya setelah pemindahan. Partisi \perintah{DEFAULT} yang berisi
banyak baris adalah tanda bahwa ada rentang tanggal yang belum Anda buatkan
partisinya.
\end{peringatan}

\langkah{B.2 Membuktikan pruning gagal}{10 menit}

Jalankan kembali query baseline apa adanya pada tabel berpartisi. Perhatikan
bahwa seluruh partisi tetap dibaca.

\begin{lstlisting}[style=sqlgaya]
EXPLAIN (ANALYZE, BUFFERS)
SELECT dp.kategori, dt.tahun_bulan, SUM(f.subtotal)
FROM   gudang.fakta_penjualan f          -- kini berpartisi
JOIN   gudang.dim_tanggal dt ON dt.tanggal_key = f.tanggal_key
JOIN   gudang.dim_produk  dp ON dp.produk_key  = f.produk_key
JOIN   gudang.dim_toko   dtk ON dtk.toko_key   = f.toko_key
WHERE  dtk.provinsi = 'Lampung'
  AND  dt.tanggal BETWEEN DATE '2024-01-01' AND DATE '2024-06-30'
GROUP  BY 1, 2;
\end{lstlisting}

Hitung berapa simpul \perintah{Seq Scan} yang muncul pada rencana. Bila
jumlahnya sama dengan jumlah partisi, pruning tidak terjadi --- penyaringan
tanggal berada pada \perintah{dim\_tanggal}, bukan pada kunci partisi.

\langkah{B.3 Membuktikan pruning berhasil}{10 menit}

Tambahkan satu predikat atas kunci partisi. Predikat ini tidak mengubah hasil
sama sekali; ia hanya memberi perencana dasar untuk memangkas.

\begin{lstlisting}[style=sqlgaya]
EXPLAIN (ANALYZE, BUFFERS)
SELECT dp.kategori, dt.tahun_bulan, SUM(f.subtotal)
FROM   gudang.fakta_penjualan f
JOIN   gudang.dim_tanggal dt ON dt.tanggal_key = f.tanggal_key
JOIN   gudang.dim_produk  dp ON dp.produk_key  = f.produk_key
JOIN   gudang.dim_toko   dtk ON dtk.toko_key   = f.toko_key
WHERE  dtk.provinsi = 'Lampung'
  AND  f.tanggal_key BETWEEN 20240101 AND 20240630   -- kunci partisi
  AND  dt.tanggal BETWEEN DATE '2024-01-01' AND DATE '2024-06-30'
GROUP  BY 1, 2;
\end{lstlisting}

Sekarang hanya enam partisi yang muncul pada rencana. Simpan kedua keluaran
\perintah{EXPLAIN} berdampingan pada \berkas{catatan-optimisasi.md} --- inilah
bukti pruning yang diminta kriteria kelulusan modul.

\begin{catatan}
Angka yang paling meyakinkan bukan waktu, melainkan jumlah partisi yang
disebut pada rencana dan total halaman yang dibaca. Bila tabel Anda memuat 36
partisi dan query hanya menyentuh 6, Anda membaca sekitar seperenam data ---
dan angka itu dapat diperiksa siapa pun tanpa menjalankan ulang query Anda.
\end{catatan}

\langkah{B.4 Menambahkan indeks pada tabel terpartisi}{5 menit}

\begin{lstlisting}[style=sqlgaya]
-- Indeks pada tabel induk otomatis dibuat pada setiap partisi.
CREATE INDEX ix_fp_toko_tanggal
  ON gudang.fakta_penjualan (toko_key, tanggal_key);

ANALYZE gudang.fakta_penjualan;
\end{lstlisting}

Ukur ulang, lalu catat sebagai tindakan kedua. Bandingkan juga ruang indeks
yang Anda bayar dengan \perintah{pg\_indexes\_size}.

\subsection{Bagian C: Membandingkan Dua Strategi Agregasi}

\langkah{C.1 Membangun aggregate table}{10 menit}

\begin{lstlisting}[style=sqlgaya]
CREATE TABLE mart.agg_penjualan_bulanan (
  tahun_bulan     INTEGER     NOT NULL,
  kategori        VARCHAR(60) NOT NULL,
  provinsi        VARCHAR(60) NOT NULL,
  unit_terjual    NUMERIC(14,2) NOT NULL,
  nilai_penjualan BIGINT      NOT NULL,
  PRIMARY KEY (tahun_bulan, kategori, provinsi)
);

INSERT INTO mart.agg_penjualan_bulanan
SELECT dt.tahun_bulan, dp.kategori, dtk.provinsi,
       SUM(f.kuantitas), SUM(f.subtotal)
FROM   gudang.fakta_penjualan f
JOIN   gudang.dim_tanggal dt ON dt.tanggal_key = f.tanggal_key
JOIN   gudang.dim_produk  dp ON dp.produk_key  = f.produk_key
JOIN   gudang.dim_toko   dtk ON dtk.toko_key   = f.toko_key
GROUP  BY 1, 2, 3;
\end{lstlisting}

Tulis juga skrip penyegaran \emph{sebagian}, yang hanya menghitung ulang satu
bulan. Kemampuan inilah yang membedakannya dari materialized view.

\begin{lstlisting}[style=sqlgaya]
DELETE FROM mart.agg_penjualan_bulanan WHERE tahun_bulan = 202406;
INSERT INTO mart.agg_penjualan_bulanan
SELECT ... WHERE f.tanggal_key BETWEEN 20240601 AND 20240630 ...;
\end{lstlisting}

\langkah{C.2 Membangun materialized view}{10 menit}

\begin{lstlisting}[style=sqlgaya]
CREATE MATERIALIZED VIEW mart.mv_penjualan_bulanan AS
SELECT dt.tahun_bulan, dp.kategori, dtk.provinsi,
       SUM(f.kuantitas) AS unit_terjual,
       SUM(f.subtotal)  AS nilai_penjualan
FROM   gudang.fakta_penjualan f
JOIN   gudang.dim_tanggal dt ON dt.tanggal_key = f.tanggal_key
JOIN   gudang.dim_produk  dp ON dp.produk_key  = f.produk_key
JOIN   gudang.dim_toko   dtk ON dtk.toko_key   = f.toko_key
GROUP  BY 1, 2, 3;

CREATE UNIQUE INDEX ux_mv_penjualan_bulanan
  ON mart.mv_penjualan_bulanan (tahun_bulan, kategori, provinsi);
\end{lstlisting}

\langkah{C.3 Membandingkan keduanya}{10 menit}

Ukur ketiga hal berikut untuk masing-masing, lalu lengkapi
Tabel~\ref{tab:m6-perbandingan}.

\begin{lstlisting}[style=sqlgaya]
-- Ruang.
SELECT pg_size_pretty(pg_total_relation_size('mart.agg_penjualan_bulanan')),
       pg_size_pretty(pg_total_relation_size('mart.mv_penjualan_bulanan'));

-- Waktu penyegaran. Bandingkan penyegaran satu bulan pada aggregate table
-- dengan penyegaran penuh materialized view.
\timing on
REFRESH MATERIALIZED VIEW mart.mv_penjualan_bulanan;
\timing off

-- Halaman yang dibaca saat menjawab pertanyaan yang sama.
EXPLAIN (ANALYZE, BUFFERS)
SELECT kategori, tahun_bulan, nilai_penjualan
FROM   mart.agg_penjualan_bulanan
WHERE  provinsi = 'Lampung' AND tahun_bulan BETWEEN 202401 AND 202406;
\end{lstlisting}

\begin{table}[htbp]
\centering
\small
\caption{Kerangka perbandingan empat tindakan optimisasi}
\label{tab:m6-perbandingan}
\begin{tabular}{@{}p{0.24\textwidth}rrrr@{}}
\toprule
\textbf{Tindakan} & \textbf{Halaman} & \textbf{Waktu 1} & \textbf{Waktu 3} &
  \textbf{Ruang} \\
\midrule
Baseline & \ldots & \ldots & \ldots & \ldots \\
Partisi bulanan & \ldots & \ldots & \ldots & \ldots \\
Partisi + indeks & \ldots & \ldots & \ldots & \ldots \\
Aggregate table & \ldots & \ldots & \ldots & \ldots \\
Materialized view & \ldots & \ldots & \ldots & \ldots \\
\bottomrule
\end{tabular}
\end{table}

\subsection{Bagian D: Menjalankan Star Query Northwind}

\langkah{D.1 Menulis query setara pada skema ternormalisasi}{10 menit}

Northwind adalah basis data operasional yang ternormalisasi --- bentuk yang
sama dengan sumber NusaMart sebelum dimodelkan secara dimensional. Tulis satu
query yang menjawab pertanyaan analitik sederhana: nilai penjualan per kategori
produk per tahun.

\begin{lstlisting}[style=sqlgaya]
EXPLAIN (ANALYZE, BUFFERS)
SELECT c.category_name,
       EXTRACT(YEAR FROM o.order_date) AS tahun,
       SUM(od.unit_price * od.quantity * (1 - od.discount)) AS nilai
FROM       order_details od
JOIN       orders     o ON o.order_id    = od.order_id
JOIN       products   p ON p.product_id  = od.product_id
JOIN       categories c ON c.category_id = p.category_id
GROUP  BY 1, 2
ORDER  BY 1, 2;
\end{lstlisting}

\langkah{D.2 Membandingkan bentuk, bukan hanya biaya}{5 menit}

Bandingkan query di atas dengan star query setara pada NusaMart. Catat pada
\berkas{catatan-optimisasi.md}:

\begin{enumerate}[leftmargin=1.4em]
  \item berapa join yang diperlukan masing-masing;
  \item berapa dalam rantai join sampai mencapai atribut kategori;
  \item ke mana perhitungan nilai penjualan diletakkan --- dihitung saat query,
        atau sudah tersimpan sebagai measure.
\end{enumerate}

\begin{catatan}
Northwind berukuran kecil, sehingga perbandingan waktu tidak bermakna dan
memang bukan itu yang dicari. Yang dibandingkan adalah \emph{bentuk} query:
berapa banyak yang harus diketahui penulisnya tentang struktur basis data
sebelum ia dapat menjawab satu pertanyaan sederhana. Inilah alasan pemodelan
dimensional ada, dan alasan itu tidak terlihat dari waktu eksekusi.
\end{catatan}

\begin{luaran}
\berkas{modul-06-optimisasi/catatan-optimisasi.md} berisi tabel
sebelum--sesudah untuk empat tindakan dalam satuan halaman dan waktu, dua
keluaran \perintah{EXPLAIN} berdampingan sebagai bukti pruning, serta
perbandingan bentuk query Northwind.
\end{luaran}

\begin{checkpoint}
Asisten menjalankan \berkas{sql/modul-06/check.sql} dan memeriksa butir berikut:
\begin{ceklis}
  \item tabel fakta sudah berpartisi menurut rentang, dengan partisi bulanan
        dan satu partisi \perintah{DEFAULT};
  \item jumlah baris sesudah pemindahan sama persis dengan sebelumnya;
  \item dataset yang dipakai berukuran sedang, bukan kecil;
  \item aggregate table dan materialized view ada pada schema \perintah{mart}
        dan keduanya menghasilkan angka yang sama;
  \item \berkas{catatan-optimisasi.md} memuat bukti pruning berupa dua keluaran
        \perintah{EXPLAIN} yang dibandingkan;
  \item pengukuran menyertakan \perintah{BUFFERS} dan menyebutkan pengaruh
        cache pada eksekusi kedua;
  \item setiap tindakan disertai keterangan apa yang dibeli dan apa yang
        dibayar.
\end{ceklis}
\end{checkpoint}

\section{Latihan Mandiri di Kelas}

\latihan{Memprediksi tindakan terbaik untuk pola filter berbeda}

Untuk setiap pola berikut, prediksi tindakan mana --- partisi, indeks,
aggregate table, atau tidak sama sekali --- yang paling menguntungkan.
Tuliskan prediksi \emph{sebelum} mengujinya, lalu uji dan jelaskan selisihnya.

\begin{enumerate}[leftmargin=1.4em]
  \item Laporan bulanan yang selalu meminta satu bulan terakhir.
  \item Pencarian satu nomor struk tertentu.
  \item Dashboard yang menampilkan total tiga tahun terakhir per kategori.
  \item Analisis yang menyaring satu produk sepanjang seluruh riwayat.
\end{enumerate}

\latihan{Menakar biaya penyegaran}

Ukur waktu penyegaran materialized view sekarang, lalu jawab:

\begin{enumerate}[leftmargin=1.4em]
  \item Bila data tumbuh dua kali lipat, kira-kira menjadi berapa lama?
  \item Bila hanya data satu hari yang berubah, berapa bagian dari pekerjaan
        penyegaran itu yang sebenarnya sia-sia?
  \item Pada ukuran data berapa penyegaran penuh menjadi tidak masuk akal
        untuk dijalankan setiap malam? Nyatakan asumsi Anda.
\end{enumerate}

\section{Tugas Individu}

\subsection{Pekerjaan Wajib}
Pilih aggregate table atau materialized view untuk beban kerja NusaMart dan
dukung pilihan dengan biaya ruang, refresh, dan pola pemakaian.

Kumpulkan \berkas{pilihan-agregasi.md} yang memuat:

\begin{enumerate}[leftmargin=1.4em]
  \item asumsi beban kerja yang dinyatakan eksplisit --- berapa kali sehari
        pertanyaan itu diajukan, dan seberapa segar data yang dituntut;
  \item pengukuran ruang keduanya, beserta perbandingannya terhadap ukuran
        tabel fakta;
  \item pengukuran waktu penyegaran penuh lawan penyegaran sebagian, beserta
        proyeksinya bila data tumbuh lima kali lipat;
  \item pilihan Anda beserta alasannya, dan --- ini yang dinilai ---
        \textbf{kondisi yang akan membuat Anda mengubah pilihan itu}.
\end{enumerate}

\begin{catatan}
Butir keempat memisahkan jawaban yang dihafal dari jawaban yang dipahami.
Tidak ada pilihan yang benar untuk semua keadaan: materialized view menang
ketika datanya kecil dan kodenya ingin ringkas; aggregate table menang ketika
data terus tumbuh dan hanya sebagian kecil yang berubah setiap hari.
Sebutkan ambang yang membuat Anda berpindah.
\end{catatan}

\subsection{Pertanyaan Analisis}

\begin{enumerate}[leftmargin=1.4em]
  \item Query Anda pada Bagian B.2 menghasilkan angka yang benar tetapi membaca
        seluruh partisi. Jelaskan mengapa kesalahan semacam ini sulit
        ditemukan di lingkungan sebenarnya, dan usulkan satu cara agar ia
        ketahuan lebih awal.
  \item Partisi bulanan memberi keuntungan bagi query yang menyaring rentang
        tanggal. Sebutkan satu pola query pada NusaMart yang justru
        \emph{dirugikan} oleh partisi bulanan, beserta alasannya.
  \item Seorang rekan melaporkan bahwa query-nya menjadi sepuluh kali lebih
        cepat setelah ia membuat indeks, dengan bukti waktu eksekusi sebelum
        dan sesudah. Sebutkan dua sebab lain yang mungkin menjelaskan
        percepatan itu, dan pengukuran apa yang akan Anda minta untuk
        memastikannya.
\end{enumerate}

\section{Format Pengumpulan dan Checklist}

Kumpulkan catatan optimisasi dengan kondisi uji, plan, halaman buffer, waktu,
dan interpretasi untuk setiap tindakan.

Seluruh berkas diletakkan pada \berkas{modul-06-optimisasi/}.

\begin{ceklis}
  \item \berkas{pre-lab.md} dan \berkas{baseline.md}
  \item \berkas{ddl-partisi.sql}
  \item \berkas{ddl-agregat.sql} --- aggregate table, skrip penyegaran
        sebagian, dan materialized view
  \item \berkas{catatan-optimisasi.md} --- tabel empat tindakan, bukti pruning,
        dan perbandingan Northwind
  \item \berkas{pilihan-agregasi.md}
  \item \berkas{jawaban-analisis.md}
  \item Setiap angka disertai keterangan kondisi ujinya: ukuran dataset,
        eksekusi keberapa, dan apakah \perintah{ANALYZE} sudah dijalankan.
\end{ceklis}

\section{Troubleshooting}

\begin{table}[htbp]
\centering
\small
\caption{Masalah yang lazim muncul pada Modul 6 dan penanganannya}
\label{tab:m6-troubleshooting}
\begin{tabular}{@{}p{0.30\textwidth}>{\raggedright\arraybackslash}p{0.62\textwidth}@{}}
\toprule
\textbf{Gejala} & \textbf{Penanganan} \\
\midrule
\perintah{no partition of relation found for row} &
  Ada baris yang tanggalnya di luar seluruh rentang partisi --- paling sering
  baris unknown berkunci $-1$ dari Modul 5. Buat partisi
  \perintah{DEFAULT}. \\
Seluruh partisi tetap dibaca &
  Penyaringan hanya ada pada \perintah{dim\_tanggal}. Tambahkan predikat atas
  \perintah{f.tanggal\_key}; ini bukan kekurangan partisi melainkan syaratnya. \\
\perintah{cannot create index on partitioned table ... CONCURRENTLY} &
  Indeks pada tabel induk tidak dapat dibuat secara concurrent. Buat biasa,
  atau buat per partisi. \\
Partisi \perintah{DEFAULT} berisi banyak baris &
  Ada rentang bulan yang belum dibuatkan partisinya. Periksa
  \perintah{MIN} dan \perintah{MAX} \perintah{tanggal\_key}, lalu lengkapi. \\
Waktu justru bertambah setelah partisi &
  Query menyentuh seluruh partisi, sehingga menanggung ongkos perencanaan
  tambahan tanpa memperoleh pruning. Periksa predikatnya. \\
Perencana memilih \perintah{Seq Scan} meski indeks ada &
  Bila query memang membaca sebagian besar tabel, ini pilihan yang benar.
  Periksa apakah selektivitas predikat Anda memang tinggi. \\
Angka aggregate table dan materialized view berbeda &
  Salah satunya belum disegarkan setelah fakta berubah. Inilah risiko yang
  disebut pada Konsep Dasar --- catat kejadiannya sebagai temuan. \\
\perintah{REFRESH ... CONCURRENTLY} ditolak &
  Materialized view belum memiliki indeks \perintah{UNIQUE}. Buat lebih
  dahulu. \\
Waktu eksekusi berubah-ubah antarpengulangan &
  Wajar. Bandingkan halaman, bukan detik, dan catat eksekusi keberapa setiap
  angka diambil. \\
\bottomrule
\end{tabular}
\end{table}

\section{Ringkasan}

Modul ini menuntut satu hal yang tidak dituntut modul sebelumnya: setiap klaim
harus disertai bukti dalam satuan yang tidak berubah-ubah. Tiga gagasan perlu
dibawa ke Modul 7.

\emph{Pertama}, partisi tidak mempercepat apa pun dengan sendirinya. Ia hanya
membuka kemungkinan untuk tidak membaca --- dan kemungkinan itu baru terpakai
bila penyaringan menyentuh kunci partisi. Query yang benar hasilnya tetapi
tidak memangkas partisi adalah kegagalan yang tidak menimbulkan galat apa pun.

\emph{Kedua}, agregasi yang disimpan lebih dahulu selalu menukar tiga hal:
ruang, waktu penyegaran, dan kesegaran data. Pilihan antara aggregate table dan
materialized view bukan soal mana yang lebih baik, melainkan soal seberapa
besar bagian data yang berubah setiap hari.

\emph{Ketiga}, cache membuat pengukuran waktu tidak dapat dipercaya. Halaman
yang disentuh tidak berubah oleh cache maupun beban laptop, dan hanya angka
semacam itu yang layak dipakai untuk membandingkan dua rancangan.

Perbandingan dengan Northwind menutup paruh perancangan praktikum ini. Sampai
sini, seluruh data berasal dari \emph{satu} sumber yang sudah rapi. Modul 7
memulai persoalan yang sama sekali berbeda: menyatukan tiga sumber operasional
yang sengaja tidak sepakat satu sama lain.
